\color{unidarkgreen}{\subsection{Максимальная токовая защита ошиновки стороны НН реактора (МТЗ НН)}}
\color{black}

\begin{enumerate}[label=\arabic{section}.\arabic{subsection}.\arabic*,  align=left, labelsep=4pt, leftmargin=0pt, itemindent=57pt]

\item
Защита предназначена для резервирования действия быстродействующей защиты ошиновки НН
стороны НН УШР, а также защит выпрямительных преобразователей ТМП.
\item
Функция выполнена двухступенчатой. Измерительные цепи защиты включены на токи фаз А, В, С
линейных выводов компенсационной обмотки УШР. При необходимости защита может быть включена
на разность токов фаз (цифровая сборка токовых цепей в «треугольник»):
\begin{equation}
i'_{A} = \frac{i_{A}-i_{B}}{\sqrt3}; \; i'_{B} = \frac{i_{B}-i_{C}}{\sqrt3}; \; i'_{C} = \frac{i_{C}-i_{A}}{\sqrt3}
\end{equation}

\item
Тип подключения защиты определяется программной накладкой \textbf{«ИО МТЗ НН»}. Каждая из
ступеней имеет регулируемые значения по току и времени срабатывания. 

\begin{figure}[H]
\centering
\includegraphics[width=0.7\textwidth,height=0.7\textheight,keepaspectratio]{\fbpath/МТЗ НН/img/MTZ_NN.pdf}
\captionof{figure}{Функциональный блок МТЗ НН}\label{pict_mtznn:MTZ_NN}
\end{figure}

\small
\begin{longtable}{|>{\centering\arraybackslash}m{5.3cm}|>{\centering\arraybackslash}m{4.2cm}|>{\centering\arraybackslash}m{1cm}|>{\centering\arraybackslash}m{3.3cm}|>{\centering\arraybackslash}m{1.8cm}|}
% Первая страница: полный заголовок
\caption{Перечень уставочных параметров МТЗ НН\hfill\vspace{-0.5\baselineskip}}\label{mtz_nn:tbl1}\\
\hline
\rowcolor{gray!30}
Наименование уставки & Диапазон & Шаг & Значение по умолчанию & Единица измерения \\
\hline
\endfirsthead

% Промежуточные страницы: «Продолжение таблицы»
\caption*{\hspace{3pt}\emph{Продолжение таблицы \ref{mtz_nn:tbl1}}\hfill\ }\vspace{-0.5\baselineskip} \\
\hline
\rowcolor{gray!30}
Наименование уставки & Диапазон & Шаг & Значение по умолчанию & Единица измерения \\
\hline
\endhead

%===m>МТЗ НН|-
%===m
\end{longtable}
\normalsize

\end{enumerate}